% SPDX-License-Identifier: GPL-3.0-or-later
% Copyright (C) 2026 Alessandro Greco
%% Русское альтернативное описание EPUB.
\altfig{Сравнение трёх архитектур. (a) Tor Browser направляет через tor только собственный трафик; другие приложения на компьютере не защищены автоматически и могут подключаться напрямую. Пунктирная рамка охватывает только браузер и tor. (b) Tails загружается как live-система с USB и работает преимущественно в памяти; правила межсетевого экрана обычно блокируют прямые соединения приложений и заставляют использовать Tor. Существуют намеренные исключения, включая небезопасный браузер Unsafe Browser для captive portal, который не использует Tor. (c) Whonix разделяет приложения в Workstation VM и tor в Gateway VM, связанные изолированной виртуальной сетью. Компрометация хостовой системы или гипервизора по-прежнему опасна.}
% Схема: браузер, live-система с межсетевым экраном, две VM
\begin{tikzpicture}[font=\sffamily\scriptsize,
  box/.style={draw,rounded corners=2pt,align=center,minimum width=2.2cm,minimum height=0.7cm},
  conf/.style={draw=torviola,dashed,rounded corners=4pt,inner sep=5pt}]
  % Tor Browser
  \node[box] (tb) at (0,1.6) {Tor Browser};
  \node[box,fill=torviola!15] (t1) at (0,0.4) {tor};
  \node[box,fill=torgrigio!15] (a1) at (0,-0.8) {другие приложения};
  \draw[->] (tb) -- node[right] {SOCKS} (t1);
  \draw[->,avviso] (a1.east) -- ++(0.7,0) node[right,align=left] {возможен прямой\\доступ в сеть};
  \node[conf,fit=(tb)(t1)] (s1) {};
  \node[above=2pt of s1] {\small (a) только трафик браузера};
  % Tails
  \begin{scope}[xshift=4.4cm]
  \node[box] (tb2) at (0,1.6) {приложения};
  \node[box,fill=torviola!15] (t2) at (0,-0.8) {tor};
  \node[box,fill=avviso!10] (fw) at (0,0.4) {firewall: блокирует обход};
  \draw[->] (tb2) -- (fw);
  \draw[->] (fw) -- (t2);
  \node[conf,fit=(tb2)(fw)] (s2) {};
  \node[above=2pt of s2] {\small (b) Tails: live USB, RAM};
  \end{scope}
  % Whonix
  \begin{scope}[xshift=8.8cm]
  \node[box] (ws) at (0,1.6) {Workstation\\(приложения)};
  \node[box,fill=torviola!15] (gw) at (0,-0.8) {Gateway\\(tor)};
  \draw[->] (ws) -- node[right,align=left] {изолированная\\виртуальная сеть} (gw);
  \node[conf,fit=(ws)] {};
  \node[conf,fit=(gw)] {};
  \node[above=12pt of ws] {\small (c) Whonix: две VM};
  \end{scope}
  \foreach \n in {t1,t2,gw} \draw[->,thick,torverde!60!black] (\n.south) -- ++(0,-0.5) node[below] {сеть Tor};
\end{tikzpicture}
